``Ballistic'' is the glide initial state propagated with \emph{no}
empirical acceleration. Running the same EA estimator on it must not invent
acceleration. Two regimes:

\begin{enumerate}
\item \textbf{Clean, default IRGN} --- recovered EA and reprop RMSE
${\approx}0$ (no hallucination on noise-free null data).
\item \textbf{Noisy + train-test split} --- 100\,Hz, $\sigma=1.0\,$m,
hyperparameters chosen by hv-block CV (\code{set\_cv()}) + GN (the
Chapter~5 deployment config): does CV overfit noise into a phantom EA?
\end{enumerate}

\examplenote{\code{examples/py/ex06\_ballistic\_recovery.py}}
\exfile{examples/py/ex06_ballistic_recovery.py}{Python}
\resultfile{results/ea_ballistic_recovery_py.csv}{Null-data control: clean and noisy+CV.}

Clean returns peak $|\mathrm{EA}|=0$, reprop $=0$ (identical to C++). The
noisy + hv-block-CV case also stays near zero --- recovered
$|\mathrm{EA}|\lesssim0.01\,$m/s$^2$, reprop ${\sim}0.1$--$0.2\,$m (exact
value carries cross-binary CV variance) against the ${\sim}35\,$km
no-EA baseline (Chapter~4). The train-test split does \textbf{not}
manufacture spurious acceleration from noise.
